% Point to the main file
\documentclass[../../Main.tex]{subfiles}

% ========== Start the document ==========
\begin{document}

\section{Section 4.6 - Cryptography}

Sorry in advance to anyone reading this since the professor never really gave a proper intro to this section and I'm too lazy right now to write an actual intro

But basically this section is about methods of encrypting data

Let's start with defining our messages

\begin{center}
    $\underbrace{\text{A}}_{00}$ \; $\underbrace{\text{B}}_{01}$ \; $\underbrace{\text{C}}_{02}$ \; $\underbrace{\text{D}}_{03}$ \; $\underbrace{\text{E}}_{04}$ \\
    $\underbrace{\text{F}}_{05}$ \; $\underbrace{\text{G}}_{06}$ \; $\underbrace{\text{H}}_{07}$ \; $\underbrace{\text{I}}_{08}$ \; $\underbrace{\text{J}}_{09}$ \\
    $\underbrace{\text{K}}_{10}$ \; $\underbrace{\text{L}}_{11}$ \; $\underbrace{\text{M}}_{12}$ \; $\underbrace{\text{N}}_{13}$ \; $\underbrace{\text{O}}_{14}$ \\
    $\underbrace{\text{P}}_{15}$ \; $\underbrace{\text{Q}}_{16}$ \; $\underbrace{\text{R}}_{17}$ \; $\underbrace{\text{S}}_{18}$ \; $\underbrace{\text{T}}_{19}$ \\
    $\underbrace{\text{U}}_{20}$ \; $\underbrace{\text{V}}_{21}$ \; $\underbrace{\text{W}}_{22}$ \; $\underbrace{\text{X}}_{23}$ \; $\underbrace{\text{Y}}_{24}$ \; $\underbrace{\text{Z}}_{25}$ \\
\end{center}

\begin{definition}[Encryption]
    Given a message $P$ (plain text), like \textbf{H E L L O}

    using some method $b$ to change the message into $C$ (cyphertext)
\end{definition}

\begin{example}
    Let's use $P$ from earlier, \textbf{H E L L O}, and get $C$ by shifting the number value of each letter by $+1$

    \begin{center}
        $\underbrace{\text{H}}_{07}$ \; $\underbrace{\text{E}}_{04}$ \; $\underbrace{\text{L}}_{11}$ \; $\underbrace{\text{L}}_{11}$ \; $\underbrace{\text{O}}_{14}$ \\
        
        ${\downarrow}_{+1}$ \\

        $\underbrace{\text{I}}_{08}$ \; $\underbrace{\text{F}}_{05}$ \; $\underbrace{\text{M}}_{12}$ \; $\underbrace{\text{M}}_{12}$ \; $\underbrace{\text{P}}_{15}$ \\
    \end{center}

    or you could do 

    $$x \mapsto ax + b$$

    This is called the \textbf{Caesar Cipher}
\end{example}

\begin{definition}[Exponentiation Cycle]
    Let's define this by doing an example

    \begin{example}
        Let's call $P$ $M$ for this example ($M$ for message)

        We need to choose a prime number bigger than 26 (since there are 26 letters)

        Let's do $p = 43$

        and choose an exponent

        Let's do $e = 5$

        Remember that our encrypted message is $C$

        Let's just do the $\underbrace{\text{H}}_{07}$ in HELLO since I'm lazy

        Let's start by \textbf{encrypting} our message.

        This means that

        $p = 43$

        $e = 5$

        $M = 7$

        \begin{align*}
            C &\equiv M^e \pmod{p} \\
            C &\equiv 7^5 \pmod{43} \\
            &\equiv 16,807 \pmod{43} \\
            \Aboxed{&\equiv 37 \pmod{43}} \\
        \end{align*}

        Boom, so 37 is our encrypted message that the receiver has to decrypt.

        Now let's look at how to \textbf{decrypt} the message
 
        We KNOW $$C \equiv M^e \pmod{p}$$

        But we WANT $$C^d \equiv M \pmod{p}$$ (trying to find $d$) because we need to find $M$

        Plugging the KNOW (first equation) into the WANT (second equation), we get

        \begin{align*}
            \left(M^e\right)^d &\equiv M \pmod{p} \\
            M^{ed} &\equiv M \pmod{p} \\
            M^{ed} \cdot M^{-1} &\equiv M \cdot M^{-1} \pmod{p} \\
            M^{ed - 1} &\equiv 1 \pmod{p}
        \end{align*}

        \begin{recall}
            From 4.4, we saw that

            $$a^k \equiv 1 \pmod{p}, \; k \equiv 0 \pmod{p - 1}$$
        \end{recall}

        In this case, $k = ed - 1$, and we see that

        \begin{align*}
            ed - 1 &\equiv 0 \pmod{p - 1} \\
            \Aboxed{ed &\equiv 1 \pmod{p - 1}}
        \end{align*}

        And since we know $e$ and $p$, we can find $d$

        \begin{align*}
            ed &\equiv 1 \pmod{p - 1} \\
            5d &\equiv 1 \pmod{42} \\
            5d &\equiv 85 \pmod{42} \tag{add $2 \cdot 42$} \\
            5d &\ \equiv 5 \cdot 17 \pmod{42} \\
            \Aboxed{d &\equiv 17 \pmod{42}} \\
        \end{align*}

        \begin{note}
            We used a shortcut way of finding the inverse but you'd typically use the Euclidean Algorithm
        \end{note}

        \begin{recall}
            Also, note that we found what $d$ is $\pmod{p - 1}$ and not $\pmod{p}$. From 4.4 again, we can see

            $$a^x \equiv a^y \pmod{p} \text{ if } x \equiv y \pmod{p - 1}$$

            There's a lot of algebra to prove that we can use this value of $d$ even though it's a different mod, but I'm too lazy to prove it so just trust
        \end{recall}

        Okay so now that we know that $d \equiv 17 \pmod{42}$

        From the decrypter's standpoint, let's recap what know so far

        $p = 43$
        
        $e = 5$

        $d = 17$

        $C = 37$

        We want to find $M$

        Let's use the equation from earlier

        $$C^d \equiv M \pmod{p}$$

        So now we just solve

        $$37^{17} \equiv M \pmod{43}$$

        \begin{align*}
            37^{17} &\equiv (-6)^{17} \pmod{43} \\
            &\equiv - (6)^{17} \pmod{43} \\
        \end{align*}

        \begin{align*}
            6 \pmod{43} &= 6 \\
            6^2 \pmod{43} &= 36 \\
            6^4 \pmod{43} &= 36^2 \\
            &= 6 \\
            6^8 \pmod{43} &= 36 \\
            6^{16} \pmod{43} &= 6 \\
        \end{align*}

        \begin{align*}
            37^{17} &\equiv - (6)^{16} \cdot 6 \pmod{43} \\
            &\equiv -6 \cdot 6 \pmod{43} \\
            &\equiv -36 \pmod{43} \\
            \Aboxed{&\equiv 7 \pmod{43}} \\
        \end{align*}

        Which is what we expect!

    \end{example}
    
\end{definition}

\begin{definition}[RSA Encryption]
    An extension of what we just did is \yellowbox{RSA Encryption}, which uses \textbf{Extended Fermat's Little Theorem} instead of normal FLT

    The way it works is you choose 2 distinct primes, $p$ and $q$

    $$p = 29, \; q = 47$$

    \begin{align*}
        N = pq &= 29 \cdot 47 \\
        &= 1363
    \end{align*}

    And choose an exponent that is relatively prime of $p$ and $q$

    $$e = 5$$

    Let's say our message is $\underbrace{\text{H}}_{07}$ again

    $$M = 7$$

    To encrypt, we do

    $$C \equiv M^e \pmod{N}$$
    
\end{definition}

% If you want, include examples of RSA and simple RSA

Okay let's recap. We just went over 2 exponentiation cyphers

\begin{enumerate}
    \item
    {
        Given a prime $p$ and exponent $e$ that are relatively prime to each other (both public)

        \begin{align*}
            C &\equiv M^e \pmod{p} \\
            C^d &\equiv M \pmod{p} \\
        \end{align*}

        \begin{align*}
            M^{ed} &\equiv M \pmod{p} \\
            M^{ed - 1} &\equiv 1 \pmod{p} \\
            \implies ed &\equiv 1 \pmod{p - 1} \\
        \end{align*}
    }

    \item
    {
        Given 2 distinct primes $p, q$ and exponent $e$ where $N = pq$ and $e$ are both public

        \begin{align*}
            C &\equiv M^e \pmod{N} \\
            C^d &\equiv M \pmod{N} \\
        \end{align*}

        \begin{align*}
            M^{ed} &\equiv M \pmod{N} \\
            M^{ed - 1} &\equiv 1 \pmod{N} \\
            \implies ed &\equiv 1 \pmod{(p - 1)(q - 1)}
        \end{align*}
    }
\end{enumerate}

Let's introduce another cypher, but let's start with an example

\begin{example}

    Let's do HELLO again

    \begin{center}
        $\underbrace{\text{H}}_{07}$ $\underbrace{\text{E}}_{04}$ $\underbrace{\text{L}}_{11}$ $\underbrace{\text{L}}_{11}$ $\underbrace{\text{O}}_{14}$ 
    \end{center}

    But let's put each letter into pairs (and let's add a random letter since HELLO is 5 letters long)

    \begin{center}
        $\boxed{\underbrace{\text{H}}_{07} \underbrace{\text{E}}_{04}}$ 
        $\boxed{\underbrace{\text{L}}_{11} \underbrace{\text{L}}_{11}}$ 
        $\boxed{\underbrace{\text{O}}_{14} \underbrace{\text{Z}}_{25}}$ 
    \end{center}

    Let's now look at our first pair / block $\boxed{\underbrace{\text{H}}_{07} \underbrace{\text{E}}_{04}}$ 

    \begin{align*}
        M &= 0704 \\
        &= 704 \\
        p &= 43, q = 59 \\
        (p - 1)(q - 1) &= 42 \cdot 58 \\
        &= 2436 \\
        N &= pq = 2537 \\
        e &= 55 \\
    \end{align*}

    Let's start by encrypting $M$, that is

    \begin{align*}
        C &\equiv M^e \pmod{N} \\
        C &\equiv 704^{55} \pmod{2537} \\
    \end{align*}

    Solving for $C$...
    
    Okay I'm too lazy to write it all out but just reduce $704^{55}$ into powers of 2 and you get

    \begin{align*}
        704^{55} &\equiv 637 \pmod{2537} \\
        C &\equiv \boxed{637} \pmod{2537} \\
    \end{align*}
    
    Now let's decrypt $C$, that is

    \begin{align*}
        M \equiv C^d \pmod{N}
    \end{align*}

    To do this, we must solve

    \begin{align*}
        ed &\equiv 1 \pmod{(p - 1)(q - 1)} \\
        55d &\equiv 1 \pmod{2436}
    \end{align*}

    Now let's find $55^{-1}$

    We do this with the Euclidean Algorithm

    \begin{align*}
        2436 &= 44 \cdot 55 + 16 \\
        55 &= 3 \cdot 16 + 7 \\
        16 &= 2 \cdot 7 + 2 \\
        7 &= 3 \cdot 2 + 1 \\
    \end{align*}

    \begin{align*}
        1 &= 7 - 3 \cdot 2 \\
        2 &= 16 - 2 \cdot 7 \\
        7 &= 55 - 3 \cdot 16 \\
        16 &= 2436 - 44 \cdot 55 \\
    \end{align*}

    \begin{align*}
        1 &= 7 - 3 \cdot (16 - 2 \cdot 7) \\
        1 &= 7 \cdot 7 - 3 \cdot 16 \\
        1 &= 7 (55 - 3 \cdot 16) - 3 \cdot 16 \\
        1 &= 7 \cdot 55 - 24 \cdot 16 \\
        1 &= 7 \cdot 55 - 24 (2436 - 44 \cdot 55) \\
        1 &= \yellowbox{1063} \cdot 55 - 24 \cdot 2436 \\
        55^{-1} &\equiv 1063 \pmod{2436} \\
        d &\equiv \boxed{1063} \pmod{2436} \\
    \end{align*}

    \begin{align*}
        M &\equiv C^d \pmod{N} \\
        M &\equiv 637^{1063} \pmod{2537} \\
        \Aboxed{M &\equiv 704} \pmod{2537} \\
    \end{align*}

\end{example}

% ========== End the document ==========
\end{document}